(1) 主轴坐标中的相位矩阵取为 \(\operatorname{diag}(e^{i\delta/2},e^{-i\delta/2})\)。旋转回实验坐标得到
\[
\mathsf J=\mathsf R(\alpha)
\begin{pmatrix}e^{i\delta/2}&0\\0&e^{-i\delta/2}\end{pmatrix}
\mathsf R(\alpha)^{\mathsf T}
=c_\delta\mathsf I+i s_\delta
\begin{pmatrix}\cos2\alpha&\sin2\alpha\\
\sin2\alpha&-\cos2\alpha\end{pmatrix},
\quad c_\delta=\cos\frac\delta2,\quad s_\delta=\sin\frac\delta2.
\]
记后一实对称矩阵为 \(\mathsf M\)。直接相乘有 \(\mathsf M^2=\mathsf I\)，所以
\(\mathsf J^\dagger\mathsf J=(c_\delta^2+s_\delta^2)\mathsf I=\mathsf I\)。
利用 \(\mathsf M\mathbf e_\sigma=e^{2i\sigma\alpha}\mathbf e_{-\sigma}\)，得到
\[
\boxed{\mathsf J\mathbf e_\sigma=c_\delta\mathbf e_\sigma+
i s_\delta e^{2i\sigma\alpha}\mathbf e_{-\sigma}}.
\]
两旋向彼此正交，故同旋向与反旋向的功率比例分别为
\(\boxed{\eta_{\rm same}=\cos^2(\delta/2)}\) 与
\(\boxed{\eta_{\rm flip}=\sin^2(\delta/2)}\)，相加严格为一。
\(\delta=\pi\) 时发生完全旋向翻转，而 \(\delta=0\) 时板对偏振无作用。

(2) 将 \(\alpha=q\phi+\alpha_0\) 代入 (1) 式，出射态成为
\[
R(r)\left[c_\delta e^{i\ell\phi}\mathbf e_\sigma+
i s_\delta e^{2i\sigma\alpha_0}
e^{i(\ell+2\sigma q)\phi}\mathbf e_{-\sigma}\right].
\]
因此转换分量获得 Pancharatnam--Berry 几何相位
\(\boxed{\Phi_{\rm PB}=2\sigma(q\phi+\alpha_0)}\)，并由
\(\hat L_z=-i\hbar\partial_\phi\) 得到
\(\boxed{L_{z,\rm flip}=(\ell+2\sigma q)\hbar}\)。
光学主轴在转过 \(\pi\) 后表示同一无向轴；绕中心一周要求
\(2\pi q=n\pi\)，所以 \(\boxed{2q\in\mathbb Z}\)，这也保证转换场的角相位单值。
\(q=0\) 是空间均匀波片，不产生新的轨道角动量；非零半整数 \(q\) 可将圆偏振翻转与整数单位的轨道角动量转移耦合。中心取向缺陷处的 Jones 局域近似失效，故题设将其从孔径中排除。

(3) 未转换分量每光子的自旋和轨道投影分别为
\(\sigma\hbar,\ell\hbar\)；转换分量分别为
\(-\sigma\hbar,(\ell+2\sigma q)\hbar\)。
由于两种旋向正交，其概率权重为 \(c_\delta^2,s_\delta^2\)，从而
\[
\begin{aligned}
\langle S_z\rangle&=\sigma\hbar(c_\delta^2-s_\delta^2)
=\sigma\hbar\cos\delta,\\
\langle L_z\rangle&=\hbar(\ell+2\sigma q\,s_\delta^2),\\
\boxed{\langle J_z\rangle}&=\hbar\bigl[\ell+\sigma+
2\sigma(q-1)s_\delta^2\bigr].
\end{aligned}
\]
光板所得角动量等于入射减去出射，故每个入射光子向板转移
\(\Delta j_z=2\sigma(1-q)\hbar s_\delta^2\)。光子通量为
\(\dot N=P/(\hbar\omega)\)，因而光板承受的平均力矩是
\[
\boxed{\tau_z=\dot N\Delta j_z
=\frac{2\sigma(1-q)P}{\omega}\sin^2\frac\delta2}.
\]
单位 \(P/\omega\) 为能量，即力矩单位。对 \(q=1\)，每次自旋改变
\(-2\sigma\hbar\) 恰由轨道角动量增加 \(2\sigma\hbar\) 抵消，\(\tau_z=0\)。
对 \(q=1/2,\delta=\pi\)，转换光子的轨道角动量增加
\(\sigma\hbar\)、自旋减少 \(2\sigma\hbar\)，板每光子获得
\(\boxed{\sigma\hbar}\)，故 \(\tau_z=\sigma P/\omega\)。
符号依赖题设圆偏振与传播方向约定。

(4) 对水平输入 \(\mathbf e_x=(1,0)^{\mathsf T}\)，直接由 Jones 矩阵得
\[
\mathsf J\mathbf e_x=
\binom{c_\delta+i s_\delta\cos2\alpha}
{i s_\delta\sin2\alpha}.
\]
水平检偏器只保留第一分量，实部与虚部正交，故
\[
\boxed{\frac{I_x(\phi)}{I_0}
=\cos^2\frac\delta2+
\sin^2\frac\delta2\cos^2(2q\phi+2\alpha_0)}.
\]
当 \(\delta=\pi\) 时为 \(\cos^2(2q\phi+2\alpha_0)\)；
由于 \(2q\) 是非零整数，一周内有 \(\boxed{4|q|}\) 个亮瓣，
例如 \(q=1/2\) 有两个、\(q=1\) 有四个。
对 \(q\ne0\)，利用完整一周内 \(\cos^2\) 的平均值 \(1/2\)，
\[
\boxed{\frac{\langle I_x\rangle_\phi}{I_0}
=\cos^2\frac\delta2+\frac12\sin^2\frac\delta2
=1-\frac12\sin^2\frac\delta2}.
\]
无检偏器时应把两正交输出分量的模方相加，
\(|c_\delta+i s_\delta\cos2\alpha|^2+
|i s_\delta\sin2\alpha|^2=1\)，故总强度均匀；
亮瓣是空间变化的偏振纹理经过检偏投影的结果，并非无损光板吸收了光能。
